\documentclass[11pt]{article}

\usepackage{amsmath,amssymb}
\usepackage[margin=1in]{geometry}
\usepackage{booktabs}
\usepackage{array}
\usepackage{hyperref}

\title{Spectral Feasibility of a Frozen Branch-Local Cochain-Laplacian Hierarchy}
\author{Tomislav Rupi\'c \\ Independent Researcher}
\date{March 2026}

\begin{document}

\maketitle

\begin{abstract}
Phase VII of the frozen HAOS-IIP program tests one narrow question: whether the already frozen branch-local operator hierarchy supports a stable descriptive spectral regime across refinement. The phase does not modify the operator family, the refinement hierarchy, or the Phase V readout observable. Track A, the spectral-feasibility program, is selected as the authoritative outcome because the frozen operator diagnostics are already clean, while observable-coupling work is better treated as a later secondary study. The tested family is the periodic DK2D block cochain Laplacian $\Delta_h$ on the deterministic hierarchy $h = 1/n_{\mathrm{side}}$ with $n_{\mathrm{side}} \in \{12,24,36,48\}$. Across these levels, the first three distinct positive low-mode bands show stable power-law fits with exponents $1.972182$, $1.972182$, and $1.924244$, all with $R^2 \ge 0.999835$. The sampled spectral radius is monotone, the nullspace estimate remains fixed at $4$, a normalized low-mode density surrogate stays tightly collapsed with maximum pairwise distance $0.041667$, and a truncated short-time trace proxy is monotone for $t \in \{0.25,0.5,1.0\}$. The result is bounded: Phase VII establishes spectral feasibility for the frozen operator hierarchy. It does not establish continuum operators, geometric invariants, universal coefficients, or physical correspondence.
\end{abstract}

\section{Scope and Authority}

This paper freezes the authority state of Phase VII in the current HAOS-IIP repository. Phase VII is constrained by two earlier freezes. Phase V defines the frozen recovery-statistics observable, and Phase VI defines the frozen branch-local operator family. Phase VII is not allowed to alter either of them. Its purpose is only to determine whether the operator hierarchy supports a stable descriptive spectral program.

The phase offers two possible tracks: a spectral-feasibility track and an observable-coupling track. The authoritative outcome in this paper is Track A, the spectral-feasibility program. Track B is not executed as an authority path in this bundle. The reason is simple: the frozen operator diagnostics already show stable nullspace behavior, monotone envelope behavior, and robust low-mode ordering, while Phase V contains a bounded robustness edge-case classifier flip that makes coupling the less clean primary statement.

\section{Frozen Inputs}

The Phase VII authority bundle references the following frozen artifacts:
\begin{itemize}
\item operator manifest: \url{phase6-operator/phase6_operator_manifest.json};
\item Phase V authority manifest: \url{phase5-readout/phase5_authoritative_manifest.json};
\item Phase VII summary: \url{phase7-spectral/phase7_spectral_summary.md};
\item Phase VII manifest: \url{phase7-spectral/phase7_spectral_manifest.json}.
\end{itemize}

The tested operator family is fixed:
\[
O_h = \Delta_h,
\]
where $\Delta_h$ is the branch-local periodic DK2D block cochain Laplacian already frozen in Phase VI. The refinement hierarchy is also fixed:
\[
h = \frac{1}{n_{\mathrm{side}}}, \qquad n_{\mathrm{side}} \in \{12,24,36,48\}.
\]
No renormalization, topology change, or hierarchy change is introduced in Phase VII.

\section{Phase VII Probe Set}

The Track A probe set consists of five deterministic checks:
\begin{enumerate}
\item low-mode scaling of the first distinct positive eigenvalue bands across refinement;
\item spectral-envelope tracking through the sampled spectral radius;
\item kernel-structure stability through the nullspace estimate;
\item normalized density-proxy collapse using an empirical CDF of the first 24 sampled positive modes after normalization by the first positive band;
\item a short-time trace proxy defined by the truncated low-mode sum
\[
\sum_i e^{-t\lambda_i}
\]
for the deterministic set $t \in \{0.25, 0.5, 1.0\}$.
\end{enumerate}

All seeds are fixed and recorded in the manifest. The sampled lowest-mode solver uses seed rule $7000 + n_{\mathrm{side}}$, and the highest-mode solver uses seed rule $9000 + n_{\mathrm{side}}$.

\section{Per-Level Diagnostic Ledger}

\begin{table}[t]
\centering
\small
\begin{tabular}{@{}rrrrrrrr@{}}
\toprule
Level & $n_{\mathrm{side}}$ & $h$ & band 1 & band 2 & band 3 & radius & nullspace \\
\midrule
R1 & 12 & 0.083333 & 0.263337 & 0.526675 & 0.982789 & 7.862310 & 4 \\
R2 & 24 & 0.041667 & 0.067853 & 0.135706 & 0.266789 & 7.965353 & 4 \\
R3 & 36 & 0.027778 & 0.030326 & 0.060652 & 0.120382 & 7.984583 & 4 \\
R4 & 48 & 0.020833 & 0.017092 & 0.034183 & 0.068074 & 7.991324 & 4 \\
\bottomrule
\end{tabular}
\caption{Per-level spectral-feasibility ledger on the frozen hierarchy.}
\label{tab:levels}
\end{table}

\begin{table}[t]
\centering
\small
\begin{tabular}{@{}lrrr@{}}
\toprule
Trace proxy & R1 & R2 & R3 \\
\midrule
$t = 0.25$ & 36.135313 & 38.939069 & 39.468497 \\
$t = 0.50$ & 32.768913 & 37.917121 & 38.947173 \\
$t = 1.00$ & 27.241831 & 35.983284 & 37.934121 \\
\bottomrule
\end{tabular}
\caption{Trace proxy values on the first three refinement levels. The R4 values continue the same monotone trend: $39.728129$, $39.458847$, and $38.927927$.}
\label{tab:trace}
\end{table}

The structural picture is clean. The first three sampled positive bands decrease monotonically under refinement, the sampled spectral radius increases mildly and monotonically, and the nullspace estimate remains fixed at $4$ at every level.

\section{Low-Mode Scaling Stability}

The first three distinct positive bands fit a power law in $h$ with excellent agreement:

\begin{table}[t]
\centering
\small
\begin{tabular}{@{}lrrr@{}}
\toprule
Band & Exponent & $R^2$ & Monotone under refinement \\
\midrule
1 & 1.972182 & 0.999979 & True \\
2 & 1.972182 & 0.999979 & True \\
3 & 1.924244 & 0.999835 & True \\
\bottomrule
\end{tabular}
\caption{Low-mode power-law fits on the frozen hierarchy.}
\label{tab:fits}
\end{table}

The exponent span across these three bands is only $0.047938$. That small spread, together with the uniformly high $R^2$ values, supports a stable descriptive low-mode scaling regime on the branch. The claim remains descriptive: Phase VII records a coherent scaling pattern across the frozen hierarchy, not a universal law.

\section{Envelope, Density Surrogate, and Trace Proxy}

The sampled spectral radius values are
\[
7.862310,\;
7.965353,\;
7.984583,\;
7.991324 ,
\]
which form a monotone refinement sequence. The nullspace estimate stays fixed at
\[
[4,4,4,4].
\]

The normalized density surrogate also stays compact after scaling. Its maximum pairwise empirical-CDF distance across all refinement pairs is $0.041667$, and the mean pairwise distance is $0.0208335$. This indicates that after normalization by the first positive band, the low-mode profile remains tightly aligned across refinement.

The short-time trace proxy is also stable in the required qualitative sense. For each tested $t$ value, the truncated trace sum increases monotonically as refinement increases. This does not establish a heat-kernel expansion or any asymptotic coefficient structure. It only shows that the branch supports a coherent trace-like diagnostic without qualitative contradictions.

\section{Bounded Interpretation}

The authority statement is narrow:
\begin{quote}
The frozen branch-local cochain-Laplacian hierarchy supports a stable descriptive spectral-feasibility regime across the declared deterministic refinement levels.
\end{quote}

This statement is supported by four combined facts:
\begin{enumerate}
\item low-mode bands fit a common scaling regime without parameter drift;
\item the sampled spectral envelope remains monotone;
\item the kernel estimate remains fixed;
\item the normalized density proxy and the trace proxy remain qualitatively stable.
\end{enumerate}

These facts are sufficient to lock a spectral-feasibility baseline for later stages. They are not sufficient to claim a continuum operator, a geometric invariant, a universal coefficient family, or any physical correspondence.

\section{Non-Claims}

This paper does not claim:
\begin{enumerate}
\item a continuum limit;
\item heat-kernel asymptotics;
\item geometric coefficients or spectral invariants;
\item quantum or relativistic correspondence;
\item observable--spectral coupling as an authority result.
\end{enumerate}

\section{Conclusion}

Phase VII closes on the spectral-feasibility track. The frozen cochain-Laplacian family shows stable low-mode scaling, stable nullspace structure, monotone spectral-envelope behavior, tight normalized low-mode collapse, and monotone short-time trace proxies across the declared refinement hierarchy. This is enough to lock the spectral-feasibility baseline for all later spectral stages on the current branch.

Phase VII establishes spectral feasibility for the frozen operator hierarchy.

\end{document}
